
\documentclass{article}
\usepackage{colortbl}
\usepackage{makecell}
\usepackage{multirow}
\usepackage{supertabular}

\begin{document}

\newcounter{utterance}

\centering \large Interaction Transcript for game `hot\_air\_balloon', experiment `air\_balloon\_survival\_en\_complexity\_easy', episode 5 with reka{-}flash{-}3.1{-}t0.0.
\vspace{24pt}

{ \footnotesize  \setcounter{utterance}{1}
\setlength{\tabcolsep}{0pt}
\begin{supertabular}{c@{$\;$}|p{.15\linewidth}@{}p{.15\linewidth}p{.15\linewidth}p{.15\linewidth}p{.15\linewidth}p{.15\linewidth}}
    \# & $\;$A & \multicolumn{4}{c}{Game Master} & $\;\:$B\\
    \hline

    \theutterance \stepcounter{utterance}  
    & & \multicolumn{4}{p{0.6\linewidth}}{
        \cellcolor[rgb]{0.9,0.9,0.9}{
            \makecell[{{p{\linewidth}}}]{
                \texttt{\tiny{[P1$\langle$GM]}}
                \texttt{You are participating in a collaborative negotiation game.} \\
\\ 
\texttt{Together with another participant, you must agree on a single set of items that will be kept. Each of you has your own view of how much each item matters to you (importance). You do not know how the other participant values the items. Additionally, you are given the effort each item demands.} \\
\texttt{You may only agree on a set if the total effort of the selected items does not exceed a shared limit:} \\
\\ 
\texttt{LIMIT = 9475} \\
\\ 
\texttt{Here are the individual item effort values:} \\
\\ 
\texttt{Item effort = \{"B21": 944, "B00": 568, "A08": 515, "A27": 239, "B94": 876, "B53": 421, "A23": 991, "C65": 281, "A89": 788, "B67": 678, "B45": 432, "C34": 409, "C12": 279, "A05": 506, "B60": 101, "C92": 686, "B69": 851, "A84": 854, "C06": 133, "B12": 192, "A16": 573, "B85": 17, "C35": 465, "B10": 772, "C15": 46, "B09": 501, "B16": 220, "C81": 404, "A58": 843, "C24": 748, "A30": 552, "A66": 843, "B75": 947, "B46": 345, "C70": 930\}} \\
\\ 
\texttt{Here is your personal view on the importance of each item:} \\
\\ 
\texttt{Item importance values = \{"B21": 138, "B00": 583, "A08": 868, "A27": 822, "B94": 783, "B53": 65, "A23": 262, "C65": 121, "A89": 508, "B67": 780, "B45": 461, "C34": 484, "C12": 668, "A05": 389, "B60": 808, "C92": 215, "B69": 97, "A84": 500, "C06": 30, "B12": 915, "A16": 856, "B85": 400, "C35": 444, "B10": 623, "C15": 781, "B09": 786, "B16": 3, "C81": 713, "A58": 457, "C24": 273, "A30": 739, "A66": 822, "B75": 235, "B46": 606, "C70": 968\}} \\
\\ 
\texttt{Goal:} \\
\\ 
\texttt{Your goal is to negotiate a shared set of items that benefits you as much as possible (i.e., maximizes total importance to YOU), while staying within the LIMIT. You are not required to make a PROPOSAL in every message {-} you can simply negotiate as well. All tactics are allowed!} \\
\\ 
\texttt{Interaction Protocol:} \\
\\ 
\texttt{You may only use the following structured formats in a message:} \\
\\ 
\texttt{PROPOSAL: \{'A', 'B', 'C', …\}} \\
\texttt{Propose keeping exactly those items.} \\
\texttt{REFUSE: \{'A', 'B', 'C', …\}} \\
\texttt{Explicitly reject opponent's proposal.} \\
\texttt{ARGUMENT: \{'...'\}} \\
\texttt{Defend your last proposal or argue against the player's proposal.} \\
\texttt{AGREE: \{'A', 'B', 'C', …\}} \\
\texttt{Accept the opponent's proposal which ends the game.} \\
\texttt{STRATEGIC REASONING: \{'...'\}} \\
\texttt{	Describe your strategic reasoning or anticipation explaining your choice of action. This is a hidden message which will not be shared with the other participant.} \\
\\ 
\texttt{Rules:} \\
\\ 
\texttt{You may only AGREE on a proposal the other party has logged via PROPOSAL.} \\
\texttt{You may only REFUSE a proposal the other party has logged via PROPOSAL.} \\
\texttt{Total effort of any PROPOSAL or AGREE set must be ≤ LIMIT.} \\
\texttt{Do NOT reveal your hidden importance scores.} \\
\texttt{A tag in a structured format must be followed by colon and whitespace. The argument must be a python set containing 0 or more strings.} \\
\texttt{So, it must be of the form TAG: \{...\}} \\
\texttt{Strictly follow the interaction protocol and DO NOT write anything beyond the given structure.} \\
\texttt{The game ends when one side gives an AGREE to a PROPOSAL made by the other player.} \\
\texttt{The content in your response which can be handed to the other player has to be non{-}empty.} \\
\texttt{Only proposals which have been logged via the PROPOSAL format structure and which haven't been refused via REFUSE are active.} \\
\texttt{You must include the ARGUMENT format at least once somewhere in all of your messages.} \\
\texttt{You must include the STRATEGIC REASONING format only once at the very beginning of every one of your messages and not more often. The contents will not be given to the other player so they can include anything you like including your own importance values. Here you should reason step by step to come up with you next move.} \\
            }
        }
    }
    & \\ \\

    \theutterance \stepcounter{utterance}  
    & \multicolumn{4}{p{0.6\linewidth}}{
        \cellcolor[rgb]{0.9,0.9,0.9}{
            \makecell[{{p{\linewidth}}}]{
                \texttt{\tiny{[P1$\rangle$GM]}}
                \texttt{<reasoning>} \\
\texttt{Let's tackle this negotiation problem. The goal is to maximize my total importance while keeping the total effort under 9475. First, I need to understand the data given.} \\
\\ 
\texttt{I have the effort for each item and my own importance values. The other player's importance isn't known, but I don't need to know it because I just want to maximize my own total importance. The challenge is selecting a subset of items that meets the effort limit and gives me the highest possible importance sum.} \\
\\ 
\texttt{First step: Calculate the total importance for each item and see which ones are high value. But since the effort is a constraint, I can't just take all high{-}importance items. Need to balance between importance and effort.} \\
\\ 
\texttt{Let me list out all items with their effort and my importance:} \\
\\ 
\texttt{Effort:} \\
\texttt{B21: 944} \\
\texttt{B00: 568} \\
\texttt{A08: 515} \\
\texttt{A27: 239} \\
\texttt{B94: 876} \\
\texttt{B53: 421} \\
\texttt{A23: 991} \\
\texttt{C65: 281} \\
\texttt{A89: 788} \\
\texttt{B67: 678} \\
\texttt{B45: 432} \\
\texttt{C34: 409} \\
\texttt{C12: 279} \\
\texttt{A05: 506} \\
\texttt{B60: 101} \\
\texttt{C92: 686} \\
\texttt{B69: 851} \\
\texttt{A84: 854} \\
\texttt{C06: 133} \\
\texttt{B12: 915} \\
\texttt{A16: 573} \\
\texttt{B85: 17} \\
\texttt{C35: 465} \\
\texttt{B10: 772} \\
\texttt{C15: 46} \\
\texttt{B09: 501} \\
\texttt{B16: 220} \\
\texttt{C81: 404} \\
\texttt{A58: 457} \\
\texttt{C24: 748} \\
\texttt{A30: 552} \\
\texttt{A66: 843} \\
\texttt{B75: 947} \\
\texttt{B46: 345} \\
\texttt{C70: 930} \\
\\ 
\texttt{Importance:} \\
\texttt{B21:138} \\
\texttt{B00:583} \\
\texttt{A08:868} \\
\texttt{A27:822} \\
\texttt{B94:783} \\
\texttt{B53:65} \\
\texttt{A23:262} \\
\texttt{C65:121} \\
\texttt{A89:508} \\
\texttt{B67:780} \\
\texttt{B45:461} \\
\texttt{C34:484} \\
\texttt{C12:668} \\
\texttt{A05:389} \\
\texttt{B60:808} \\
\texttt{C92:215} \\
\texttt{B69:97} \\
\texttt{A84:500} \\
\texttt{C06:30} \\
\texttt{B12:915} \\
\texttt{A16:856} \\
\texttt{B85:400} \\
\texttt{C15:781} \\
\texttt{B09:786} \\
\texttt{B16:3} \\
            }
        }
    }
    & & \\ \\

    \theutterance \stepcounter{utterance}  
    & & \multicolumn{4}{p{0.6\linewidth}}{
        \cellcolor[rgb]{0.9,0.9,0.9}{
            \makecell[{{p{\linewidth}}}]{
                \texttt{\tiny{[P1$\langle$GM]}}
                \texttt{Your response did not start with the proper strategic reasoning tag at the very beginning of your response. The very first structured format must be of the form STRATEGIC REASONING: \{...\}. Try again.} \\
            }
        }
    }
    & \\ \\

    \theutterance \stepcounter{utterance}  
    & \multicolumn{4}{p{0.6\linewidth}}{
        \cellcolor[rgb]{0.9,0.9,0.9}{
            \makecell[{{p{\linewidth}}}]{
                \texttt{\tiny{[P1$\rangle$GM]}}
                \texttt{<reasoning>} \\
\texttt{Let me start by understanding the problem again. The goal is to select a subset of items such that the total effort doesn't exceed 9475, and the sum of my importance values for those items is maximized. Since the other player's preferences are unknown, I don't need to consider their importance, just focus on mine.} \\
\\ 
\texttt{First, I need to calculate the importance{-}to{-}effort ratio for each item. This ratio will help me prioritize items that give me the most importance per unit effort. High ratios mean better value. Let me list out each item with its effort and importance, then compute the ratio.} \\
\\ 
\texttt{Starting with B21: effort 944, importance 138. Ratio is 138/944 ≈ 0.146. That's low. B00: 568 effort, 583 importance. 583/568 ≈ 1.026. That's better. A08: 515 effort, 868 importance. 868/515 ≈ 1.686. High ratio. A27: 239 effort, 822 importance. 822/239 ≈ 3.44. Wow, that's really good. B94: 876 effort, 783 importance. 783/876 ≈ 0.896. Lower. B53: 421 effort, 65 importance. 65/421 ≈ 0.154. Not good. A23: 991 effort, 262 importance. 262/991 ≈ 0.264. B67: 678 effort, 780 importance. 780/678 ≈ 1.149. C34: 409 effort, 484 importance. 484/409 ≈ 1.184. C12: 279 effort, 668 importance. 668/279 ≈ 2.393. High. A05: 506 effort, 389 importance. 389/506 ≈ 0.768. B60: 101 effort, 808 importance. 808/101 ≈ 8.01. That's extremely high. C92: 686 effort, 215 importance. 215/686 ≈ 0.313. B69: 851 effort, 97 importance. 97/851 ≈ 0.114. A84: 854 effort, 500 importance. 500/854 ≈ 0.586. C06: 133 effort, 30} \\
            }
        }
    }
    & & \\ \\

    \theutterance \stepcounter{utterance}  
    & & \multicolumn{4}{p{0.6\linewidth}}{
        \cellcolor[rgb]{0.9,0.9,0.9}{
            \makecell[{{p{\linewidth}}}]{
                \texttt{\tiny{[P1$\langle$GM]}}
                \texttt{Your response did not start with the proper strategic reasoning tag at the very beginning of your response. The very first structured format must be of the form STRATEGIC REASONING: \{...\}. Try again.} \\
            }
        }
    }
    & \\ \\

    \theutterance \stepcounter{utterance}  
    & \multicolumn{4}{p{0.6\linewidth}}{
        \cellcolor[rgb]{0.9,0.9,0.9}{
            \makecell[{{p{\linewidth}}}]{
                \texttt{\tiny{[P1$\rangle$GM]}}
                \texttt{<reasoning>} \\
\texttt{Let me start by understanding the problem again. The goal is to select a subset of items such that the total effort doesn't exceed 9475, and the sum of my importance values for those items is maximized. Since the other player's preferences are unknown, I don't need to consider their importance, just focus on mine.} \\
\\ 
\texttt{First, I need to calculate the importance{-}to{-}effort ratio for each item. This ratio will help me prioritize items that give me the most importance per unit effort. High ratios mean better value. Let me list out each item with its effort and importance, then compute the ratio.} \\
\\ 
\texttt{Starting with B21: effort 944, importance 138. Ratio is 138/944 ≈ 0.146. That's low. B00: 568 effort, 583 importance. 583/568 ≈ 1.026. That's better. A08: 515 effort, 868 importance. 868/515 ≈ 1.686. High ratio. A27: 239 effort, 822 importance. 822/239 ≈ 3.44. Wow, that's really good. B94: 876 effort, 783 importance. 783/876 ≈ 0.896. Lower. B53: 421 effort, 65 importance. 65/421 ≈ 0.154. Not good. A23: 991 effort, 262 importance. 262/991 ≈ 0.264. B67: 678 effort, 780 importance. 780/678 ≈ 1.149. C34: 409 effort, 484 importance. 484/409 ≈ 1.184. C12: 279 effort, 668 importance. 668/279 ≈ 2.393. High. A05: 506 effort, 389 importance. 389/506 ≈ 0.768. B60: 101 effort, 808 importance. 808/101 ≈ 8.01. That's extremely high. C92: 686 effort, 215 importance. 215/686 ≈ 0.313. B69: 851 effort, 97 importance. 97/851 ≈ 0.114. A84: 854 effort, 500 importance. 500/854 ≈ 0.586. C06: 133 effort, 30} \\
            }
        }
    }
    & & \\ \\

    \theutterance \stepcounter{utterance}  
    & & & \multicolumn{2}{p{0.3\linewidth}}{
        \cellcolor[rgb]{0.9,0.9,0.9}{
            \makecell[{{p{\linewidth}}}]{
                \texttt{\tiny{[GM$|$GM]}}
                \texttt{abort} \\
            }
        }
    }
    & & \\ \\

    \theutterance \stepcounter{utterance}  
    & & & \multicolumn{2}{p{0.3\linewidth}}{
        \cellcolor[rgb]{0.9,0.9,0.9}{
            \makecell[{{p{\linewidth}}}]{
                \texttt{\tiny{[GM$|$GM]}}
                \texttt{end game} \\
            }
        }
    }
    & & \\ \\

\end{supertabular}
}

\end{document}
